\documentclass[11pt,a4paper]{article}

\usepackage[utf8]{inputenc}
\usepackage[T1]{fontenc}
\usepackage[ngerman]{babel}

\usepackage{amsmath,amssymb,amsthm,mathtools}
\usepackage[a4paper,margin=2.5cm]{geometry}
\usepackage{lmodern}
\usepackage{microtype}
\usepackage{hyperref}
\usepackage{enumitem}

\title{Klassifikation geschichteter Radikalzerlegungen\\ und ihrer Radikalreste}
\author{Kollaborative Ausarbeitung}
\date{\today}

\newtheorem{satz}{Satz}
\newtheorem{lemma}[satz]{Lemma}
\newtheorem{proposition}[satz]{Proposition}
\newtheorem{korollar}[satz]{Korollar}
\theoremstyle{definition}
\newtheorem{definition}[satz]{Definition}
\newtheorem{beispiel}[satz]{Beispiel}
\theoremstyle{remark}
\newtheorem{bemerkung}[satz]{Bemerkung}

\newcommand{\N}{\mathbb{N}}
\newcommand{\F}{\mathbb{F}_2}
\newcommand{\rad}{\operatorname{rad}}
\newcommand{\Typ}{\operatorname{Typ}}
\newcommand{\RRest}{\mathcal{R}}
\newcommand{\vp}{v_p}
\newcommand{\ggT}{\operatorname{ggT}}
\newcommand{\epsA}{\varepsilon_A}
\newcommand{\epsB}{\varepsilon_B}
\newcommand{\epsC}{\varepsilon_C}

\begin{document}
	\maketitle
	
	\begin{abstract}
		Nach Abspaltung der Faktoren \(2\) und \(3\) besitzt jede natürliche Zahl eine eindeutige
		geschichtete Radikalzerlegung. Jede Radikalschicht ist quadratfrei und zu \(6\) teilerfremd.
		Durch Reduktion der Primteiler modulo \(12\) auf die Familien \(E,A,B,C\) und anschließende
		Paritätsreduktion erhält jede Schicht genau einen reduzierten Typ in
		\[
		\mathcal T:=\{E,A,B,C,AB,AC,BC\}.
		\]
		Die Folge dieser Typen heißt Radikalrest der Zahl. Der zentrale Klassifikationssatz besagt,
		dass die Menge aller reduzierten Radikalreste genau aus den endlichen Wörtern über dem Alphabet
		\(\mathcal T\) besteht.
	\end{abstract}
	
	\section{Geschichtete Radikalzerlegung}
	
	\begin{definition}[Abspaltung des \(2\)- und \(3\)-Anteils]
		Für \(n\in\N\) setzen wir
		\[
		\alpha:=v_2(n),\qquad \beta:=v_3(n),\qquad m:=\frac{n}{2^\alpha 3^\beta}.
		\]
		Dann gilt
		\[
		n=2^\alpha 3^\beta m
		\qquad\text{und}\qquad
		\ggT(m,6)=1.
		\]
	\end{definition}
	
	\begin{definition}[Radikal]
		Für \(x\in\N\) definieren wir
		\[
		\rad(x):=\prod_{p\mid x} p.
		\]
		Insbesondere ist \(\rad(x)\) quadratfrei.
	\end{definition}
	
	\begin{definition}[Radikalschichten]
		Sei \(m\in\N\) mit \(\ggT(m,6)=1\) und
		\[
		m=\prod_p p^{v_p(m)}.
		\]
		Für
		\[
		h:=\max_{p\mid m} v_p(m)
		\]
		definieren wir die \(k\)-te Radikalschicht durch
		\[
		R_k(m):=\prod_{\substack{p\mid m\\ v_p(m)\ge k}} p,
		\qquad k=1,\dots,h.
		\]
	\end{definition}
	
	\begin{lemma}
		Für jedes \(m\in\N\) mit \(\ggT(m,6)=1\) gilt
		\[
		m=\prod_{k=1}^{h}R_k(m).
		\]
	\end{lemma}
	
	\begin{proof}
		Eine Primzahl \(p\mid m\) erscheint genau in denjenigen Schichten \(R_k(m)\), für die
		\(v_p(m)\ge k\) gilt. Dies sind genau \(v_p(m)\) viele Schichten. Also ist der Exponent von \(p\)
		im Produkt \(\prod_{k=1}^{h}R_k(m)\) gleich \(v_p(m)\). Damit stimmen beide Seiten in ihrer
		Primfaktorzerlegung überein.
	\end{proof}
	
	\begin{lemma}
		Die Folge der Radikalschichten ist divisibilitätsabsteigend:
		\[
		R_{k+1}(m)\mid R_k(m)
		\qquad\text{für alle }k=1,\dots,h-1.
		\]
	\end{lemma}
	
	\begin{proof}
		Ist \(p\mid R_{k+1}(m)\), so gilt \(v_p(m)\ge k+1\), also insbesondere \(v_p(m)\ge k\). Daher
		folgt \(p\mid R_k(m)\). Also teilt \(R_{k+1}(m)\) die Schicht \(R_k(m)\).
	\end{proof}
	
	\begin{satz}[Eindeutige geschichtete Radikalzerlegung]
		Zu jeder natürlichen Zahl \(n\in\N\) existieren eindeutig bestimmte nichtnegative ganze Zahlen
		\(\alpha,\beta\) und eindeutig bestimmte quadratfreie Zahlen
		\[
		R_1,\dots,R_h
		\]
		mit
		\[
		n=2^\alpha 3^\beta \prod_{k=1}^{h}R_k,
		\qquad
		R_{k+1}\mid R_k,
		\qquad
		\ggT(R_k,6)=1.
		\]
	\end{satz}
	
	\begin{proof}
		Die Zahlen \(\alpha=v_2(n)\) und \(\beta=v_3(n)\) sind eindeutig bestimmt. Damit ist
		\(m=n/(2^\alpha 3^\beta)\) eindeutig gegeben. Die Radikalschichten \(R_k=R_k(m)\) sind
		definitionsgemäß eindeutig und erfüllen nach den beiden vorigen Lemmata die gewünschten
		Eigenschaften.
	\end{proof}
	
	\section{Die Familien \(E,A,B,C\) modulo \(12\)}
	
	\begin{definition}[Primfamilien]
		Für Primzahlen \(p\ge 5\) definieren wir
		\[
		E:=\{p:\ p\equiv 1 \pmod{12}\},\qquad
		A:=\{p:\ p\equiv 5 \pmod{12}\},
		\]
		\[
		B:=\{p:\ p\equiv 7 \pmod{12}\},\qquad
		C:=\{p:\ p\equiv 11 \pmod{12}\}.
		\]
	\end{definition}
	
	\begin{lemma}
		In \((\mathbb Z/12\mathbb Z)^\times\) gelten die Relationen
		\[
		A^2=B^2=C^2=E,
		\]
		\[
		AB=C,\qquad AC=B,\qquad BC=A,
		\]
		\[
		ABC=E.
		\]
	\end{lemma}
	
	\begin{proof}
		Dies folgt direkt aus
		\[
		5^2\equiv 7^2\equiv 11^2\equiv 1 \pmod{12},
		\]
		\[
		5\cdot 7\equiv 11,\qquad 5\cdot 11\equiv 7,\qquad 7\cdot 11\equiv 5 \pmod{12},
		\]
		sowie
		\[
		5\cdot 7\cdot 11\equiv 1 \pmod{12}.
		\]
	\end{proof}
	
	\begin{definition}[Paritätssignatur]
		Sei \(R\) quadratfrei mit \(\ggT(R,6)=1\). Dann setzen wir
		\[
		\epsA(R):=\#\{p\mid R:\ p\in A\}\bmod 2,
		\]
		\[
		\epsB(R):=\#\{p\mid R:\ p\in B\}\bmod 2,
		\]
		\[
		\epsC(R):=\#\{p\mid R:\ p\in C\}\bmod 2.
		\]
		Die Paritätssignatur von \(R\) ist
		\[
		\sigma(R):=(\epsA(R),\epsB(R),\epsC(R))\in\F^3.
		\]
	\end{definition}
	
	\begin{definition}[Reduzierter Typ]
		Wir definieren
		\[
		\mathcal T:=\{E,A,B,C,AB,AC,BC\}.
		\]
		Für eine quadratfreie Zahl \(R\) mit \(\ggT(R,6)=1\) sei \(\Typ(R)\in\mathcal T\) gegeben durch
		\[
		(0,0,0)\mapsto E,\quad
		(1,0,0)\mapsto A,\quad
		(0,1,0)\mapsto B,\quad
		(0,0,1)\mapsto C,
		\]
		\[
		(1,1,0)\mapsto AB,\quad
		(1,0,1)\mapsto AC,\quad
		(0,1,1)\mapsto BC,\quad
		(1,1,1)\mapsto E.
		\]
	\end{definition}
	
	\begin{proposition}[Siebentypen-Klassifikation]
		Jede quadratfreie Zahl \(R\) mit \(\ggT(R,6)=1\) besitzt genau einen reduzierten Typ
		\[
		\Typ(R)\in\mathcal T.
		\]
		Insbesondere gibt es genau sieben reduzierte Typen.
	\end{proposition}
	
	\begin{proof}
		Die Paritätssignatur \(\sigma(R)\) ist durch \(R\) eindeutig bestimmt. Formal ergeben sich acht
		binäre Muster in \(\F^3\). Die beiden Muster \((0,0,0)\) und \((1,1,1)\) werden jedoch wegen
		\(ABC=E\) demselben neutralen Typ \(E\) zugeordnet. Daher bleiben genau sieben reduzierte Typen.
	\end{proof}
	
	\section{Radikalreste}
	
	\begin{definition}[Radikalrest]
		Sei
		\[
		n=2^\alpha 3^\beta \prod_{k=1}^{h}R_k
		\]
		die geschichtete Radikalzerlegung von \(n\). Dann heißt
		\[
		\RRest(n):=\bigl(\Typ(R_1),\Typ(R_2),\dots,\Typ(R_h)\bigr)
		\]
		der \emph{Radikalrest} von \(n\).
	\end{definition}
	
	\begin{satz}[Existenz und Eindeutigkeit des Radikalrestes]
		Zu jeder natürlichen Zahl \(n\in\N\) existiert genau ein Radikalrest
		\[
		\RRest(n)\in\mathcal T^{<\infty},
		\]
		also genau ein endliches Wort über dem Alphabet \(\mathcal T\).
	\end{satz}
	
	\begin{proof}
		Die geschichtete Radikalzerlegung von \(n\) ist eindeutig. Jede Schicht besitzt nach der vorigen
		Proposition einen eindeutig bestimmten reduzierten Typ. Folglich ist die gesamte Typenfolge
		eindeutig bestimmt.
	\end{proof}
	
	\section{Klassifikation aller Radikalreste}
	
	\begin{definition}[Endliche Wörter]
		Für eine Menge \(\mathcal T\) bezeichne
		\[
		\mathcal T^{<\infty}
		\]
		die Menge aller endlichen Wörter über \(\mathcal T\), also aller endlichen Folgen
		\[
		(t_1,\dots,t_h),\qquad h\ge 1,\quad t_k\in\mathcal T.
		\]
	\end{definition}
	
	\begin{bemerkung}
		Man kann auf Wunsch auch das leere Wort zulassen und ihm die Zahl \(1\) zuordnen. Im vorliegenden
		Text betrachten wir für \(n\ge 2\) nur nichtleere Radikalreste.
	\end{bemerkung}
	
	\begin{satz}[Klassifikation der reduzierten Radikalreste]
		Die Menge aller reduzierten Radikalreste natürlicher Zahlen ist genau
		\[
		\mathcal T^{<\infty}.
		\]
		Das heißt:
		\begin{enumerate}[label=\textup{(\roman*)}]
			\item Zu jeder natürlichen Zahl \(n\) ist \(\RRest(n)\) ein endliches Wort über \(\mathcal T\).
			\item Umgekehrt wird jedes endliche Wort über \(\mathcal T\) als Radikalrest einer geeigneten
			natürlichen Zahl realisiert.
		\end{enumerate}
	\end{satz}
	
	\begin{proof}
		Aussage \textup{(i)} ist bereits gezeigt.
		
		Für \textup{(ii)} sei
		\[
		(t_1,\dots,t_h)\in\mathcal T^{<\infty}
		\]
		ein beliebiges Wort. Wir wählen eine kanonische Hebung
		\[
		\lambda:\mathcal T\to\F^3
		\]
		durch
		\[
		\lambda(E)=000,\quad \lambda(A)=100,\quad \lambda(B)=010,\quad \lambda(C)=001,
		\]
		\[
		\lambda(AB)=110,\quad \lambda(AC)=101,\quad \lambda(BC)=011.
		\]
		Setze
		\[
		s_k:=\lambda(t_k)\in\F^3
		\qquad (k=1,\dots,h),
		\qquad
		s_{h+1}:=000.
		\]
		Definiere ferner
		\[
		q_k:=s_k+s_{k+1}\in\F^3
		\qquad (k=1,\dots,h),
		\]
		wobei die Addition in \(\F^3\) komponentenweise modulo \(2\) erfolgt.
		
		Nun wählen wir für jedes \(k\) eine quadratfreie Zahl \(Q_k\) mit \(\ggT(Q_k,6)=1\), deren
		Paritätssignatur gerade \(q_k\) ist. Dies ist stets möglich. Beispielsweise kann man nehmen:
		\[
		000\rightsquigarrow 13,\quad
		100\rightsquigarrow 5,\quad
		010\rightsquigarrow 7,\quad
		001\rightsquigarrow 11,
		\]
		\[
		110\rightsquigarrow 5\cdot 7,\quad
		101\rightsquigarrow 5\cdot 11,\quad
		011\rightsquigarrow 7\cdot 11.
		\]
		Setze nun
		\[
		m:=\prod_{k=1}^{h}Q_k^k.
		\]
		Dann ist
		\[
		R_k(m)=\prod_{j\ge k}Q_j.
		\]
		Die Paritätssignatur von \(R_k(m)\) ist daher
		\[
		q_k+q_{k+1}+\cdots+q_h
		=(s_k+s_{k+1})+(s_{k+1}+s_{k+2})+\cdots+(s_h+s_{h+1})
		=s_k+s_{h+1}=s_k,
		\]
		da sich alle Zwischenterme in \(\F^3\) paarweise aufheben und \(s_{h+1}=000\) ist. Also gilt
		\[
		\Typ(R_k(m))=t_k
		\qquad\text{für alle }k=1,\dots,h.
		\]
		Somit ist
		\[
		\RRest(m)=(t_1,\dots,t_h),
		\]
		und das Wort wird realisiert.
	\end{proof}
	
	\begin{korollar}
		Es gibt keine weiteren verborgenen Konsistenzbedingungen für reduzierte Radikalreste: Jede endliche
		Folge von Symbolen aus
		\[
		\{E,A,B,C,AB,AC,BC\}
		\]
		tritt als Radikalrest mindestens einer natürlichen Zahl auf.
	\end{korollar}
	
	\section{Beispiele}
	
	\begin{beispiel}
		Für
		\[
		n=2^2\cdot 3\cdot 5^3\cdot 7^2\cdot 13
		\]
		ist
		\[
		m=\frac{n}{2^2\cdot 3}=5^3\cdot 7^2\cdot 13,
		\]
		also
		\[
		R_1=5\cdot 7\cdot 13,\qquad R_2=5\cdot 7,\qquad R_3=5.
		\]
		Daraus folgt
		\[
		\Typ(R_1)=AB,\qquad \Typ(R_2)=AB,\qquad \Typ(R_3)=A,
		\]
		also
		\[
		\RRest(n)=(AB,AB,A).
		\]
	\end{beispiel}
	
	\begin{beispiel}
		Das Wort
		\[
		(A,AB,E,C)
		\]
		ist realisierbar. Wähle dazu
		\[
		s_1=100,\quad s_2=110,\quad s_3=000,\quad s_4=001,\quad s_5=000.
		\]
		Dann erhält man
		\[
		q_1=010,\quad q_2=110,\quad q_3=001,\quad q_4=001.
		\]
		Man kann also wählen
		\[
		Q_1=7,\quad Q_2=5\cdot 7,\quad Q_3=11,\quad Q_4=11.
		\]
		Damit ist
		\[
		m=Q_1^1Q_2^2Q_3^3Q_4^4
		\]
		eine Zahl mit
		\[
		\RRest(m)=(A,AB,E,C).
		\]
	\end{beispiel}
	
	\begin{beispiel}
		Für
		\[
		n=5\cdot 7\cdot 11
		\]
		gibt es nur eine Schicht
		\[
		R_1=5\cdot 7\cdot 11.
		\]
		Das formale Muster ist \(ABC\), reduziert sich aber zu \(E\). Also
		\[
		\RRest(n)=(E).
		\]
	\end{beispiel}
	
	\section{Schluss}
	
	Die hier dargestellte Theorie trennt vier Ebenen:
	\begin{enumerate}[label=\arabic*.]
		\item die gewöhnliche Primfaktorzerlegung,
		\item die geschichtete Radikalzerlegung,
		\item die Reduktion jeder Schicht auf einen der sieben Typen,
		\item die Klassifikation aller Radikalreste als freie Wörter über dem Sieben-Typen-Alphabet.
	\end{enumerate}
	
	Damit erhält jede natürliche Zahl eine kanonische endliche Signaturfolge
	\[
	\RRest(n)\in\{E,A,B,C,AB,AC,BC\}^{<\infty},
	\]
	und umgekehrt kommt jede solche Folge tatsächlich vor.
	
\end{document}